\documentclass[11pt,a4paper]{article}
\usepackage{fontspec,amsmath,amssymb,amsthm,microtype}
\setmainfont{lmroman10-regular.otf}[BoldFont=lmroman10-bold.otf,ItalicFont=lmroman10-italic.otf,BoldItalicFont=lmroman10-bolditalic.otf]
\usepackage{xcolor}
\usepackage[most]{tcolorbox}
\usepackage{mathrsfs}
\providecommand{\lt}{<}
\providecommand{\gt}{>}
\setmonofont{DejaVuSansMono.ttf}[BoldFont=DejaVuSansMono-Bold.ttf,ItalicFont=DejaVuSansMono-Oblique.ttf,BoldItalicFont=DejaVuSansMono-BoldOblique.ttf,Scale=0.86]
\newfontfamily\leanfallback{STIXTwoMath-Regular.otf}[Scale=0.86]
\definecolor{leanrule}{gray}{0.62}
\definecolor{leanback}{gray}{0.962}
\definecolor{leantitle}{gray}{0.22}
\definecolor{leanmuted}{gray}{0.38}
\newtcolorbox{leanbox}[1]{enhanced,breakable,colback=leanback,colframe=leanrule,
  boxrule=0.4pt,arc=0.6mm,left=2.4mm,right=2.4mm,top=1.8mm,bottom=1.8mm,
  before skip=0.8em,after skip=1em,fonttitle=\small\bfseries,coltitle=leantitle,
  colbacktitle=leanback,titlerule=0.3pt,title={#1}}
\newlength{\leancharwidth}
\newenvironment{leancode}{\par\smallskip\begingroup\ttfamily\footnotesize\raggedright
  \settowidth{\leancharwidth}{M}\setlength{\parskip}{0pt}\setlength{\parindent}{0pt}
  \frenchspacing}{\par\endgroup\smallskip}
\newcommand{\leanline}[2]{\hangindent=\dimexpr#1\leancharwidth+4\leancharwidth\relax
  \hangafter=1\noindent\hspace*{#1\leancharwidth}#2\par}
\newcommand{\leansource}[1]{\par{\footnotesize\color{leanmuted}#1}\par}
\newcommand{\leanlabel}[1]{\par\smallskip\noindent{\small\bfseries #1}\par}
\usepackage[margin=25mm]{geometry}
\usepackage[unicode,hidelinks]{hyperref}
\hypersetup{pdfauthor={LangevinLean Project}}
\setlength{\emergencystretch}{3em}
\setlength{\parskip}{0.4em}
\setlength{\parindent}{0pt}
\allowdisplaybreaks[2]
\raggedbottom
\title{\texorpdfstring{Convergence in a Hilbert space of random variables}{Convergence in a Hilbert space of random variables}}
\author{LangevinLean Project}
\date{}
\begin{document}\maketitle
Draft mathematical exposition; these notes may contain mistakes.


Let \(\mathcal H=L^2(P;H)\) and let the gradient act by pointwise composition. Supply continuous reference segments \(Y_{k,t}\), joining at grid endpoints, whose grid laws equal \(\pi\). Their endpoint noise is the innovation driving the numerical step. The reference integral equation and the bounds \(\|g(Y_{k,t})\|_2\le\sqrt{Ld}\) and \(\|B_{k,t}\|_2\le\sqrt{2d}\sqrt t\) give a defect norm at most \((33/20)Lh\sqrt{hd}\) when \(Lh\le2\).



The common-noise subtraction and the deterministic Euler contraction give an affine recurrence in \(\mathcal H\). Its geometric sum bounds the transport distance of the numerical law from \(\pi\). The two step ranges use factors \(1-mh\) and \(Lh-1\), respectively. The reference and its stationary laws are hypotheses here; the Brownian construction later supplies them.


\section{\texorpdfstring{The recursion on square-integrable random variables}{The recursion on square-integrable random variables}}


For an initial \(X_0\in L^2(P;H)\) and innovations \(\xi_k\in L^2(P;H)\), set \(X_{k+1}=X_k-h\widehat{\nabla U}(X_k)+\sqrt{2h}\,\xi_k\).



The initial iterate is the supplied random vector.



Each successor uses the Euler update using the given gradient.



The drift lift evaluates the original gradient almost everywhere, so this recursion represents the ordinary ULA method on random inputs. The reference equation and defect integral are Bochner equations in this same Hilbert space. Each step stays square-integrable by the Lipschitz growth estimate for the gradient.

\begin{leanbox}{Lean: \texttt{lpChain}}
\leansource{\texttt{lpChain} \quad UnadjustedLangevin/\allowbreak{}HilbertDalalyan.lean, lines 34–38}
\begin{leancode}
\leanline{0}{def\ lpChain\ (U\ :\ H\ →\ ℝ)\ (h\ :\ ℝ)\ \{L\ :\ ℝ≥0\}}
\leanline{4}{(hLip\ :\ LipschitzWith\ L\ (gradient\ U))\ (initial\ :\ Lp\ H\ 2\ P)}
\leanline{4}{(noise\ :\ ℕ\ →\ Lp\ H\ 2\ P)\ :\ ℕ\ →\ Lp\ H\ 2\ P}
\leanline{2}{|\ 0\ =>\ initial}
\leanline{2}{|\ n\ +\ 1\ =>\ lpUlaUpdate\ P\ U\ h\ hLip\ (lpChain\ U\ h\ hLip\ initial\ noise\ n)\ (noise\ n)}
\end{leancode}
\leanlabel{How the Lean proof follows the argument}
\texttt{lp\allowbreak{}Chain} is a recursion in the Hilbert space \texttt{Lp H 2 P}. The zero branch returns \texttt{initial}; the successor applies \texttt{lp\allowbreak{}Ula\allowbreak{}Update} to the preceding state and \texttt{noise n}. The proof \texttt{h\allowbreak{}Lip} is passed to that update so composition with the drift remains square integrable. Gaussianity is not required to define this recursion.
\end{leanbox}

\section{\texorpdfstring{The successor update}{The successor update}}


The next iterate is exactly one \(L^2\) ULA update of the current iterate.



Evaluate the successor clause of the definition.



The successor step uses the innovation at the same index as the reference endpoint. Subtracting a reference equation using that same innovation removes the stochastic increment exactly. The resulting error equation therefore contains only the deterministic Euler difference and the local defect.

\begin{leanbox}{Lean: \texttt{lpChain\_\allowbreak{}succ}}
\leansource{\texttt{lpChain\_\allowbreak{}succ} \quad UnadjustedLangevin/\allowbreak{}HilbertDalalyan.lean, lines 41–45}
\begin{leancode}
\leanline{0}{@[simp]\ theorem\ lpChain\_succ\ (U\ :\ H\ →\ ℝ)\ (h\ :\ ℝ)\ \{L\ :\ ℝ≥0\}}
\leanline{4}{(hLip\ :\ LipschitzWith\ L\ (gradient\ U))\ (initial\ :\ Lp\ H\ 2\ P)}
\leanline{4}{(noise\ :\ ℕ\ →\ Lp\ H\ 2\ P)\ (n\ :\ ℕ)\ :}
\leanline{4}{lpChain\ P\ U\ h\ hLip\ initial\ noise\ (n\ +\ 1)\ =}
\leanline{6}{lpUlaUpdate\ P\ U\ h\ hLip\ (lpChain\ P\ U\ h\ hLip\ initial\ noise\ n)\ (noise\ n)}
\end{leancode}
\leanlabel{How the Lean proof follows the argument}
\texttt{lp\allowbreak{}Chain\_\allowbreak{}succ} exposes the successor branch of \texttt{lp\allowbreak{}Chain}. Its proof is \texttt{rfl}, because both sides reduce to exactly the same expression after evaluating the recursive definition.
\end{leanbox}

\section{\texorpdfstring{The continuous reference assumptions}{The continuous reference assumptions}}


For an \(L\)-Lipschitz gradient, step \(h\), noise vectors \(\xi_k\in\mathcal H\), and a measure \(\pi\), a reference consists of segments \(Y_{k,t}\in\mathcal H\) and increments \(B_{k,t}\in\mathcal H\) satisfying the following conditions.



Supply the two families \((Y_{k,t})\) and \((B_{k,t})\) in the Hilbert space \(\mathcal H\).



Require continuity of \(t\mapsto Y_{k,t}\) in the \(L^2\) norm for every \(k\).



Require compatible endpoints, common endpoint noise, and stationary grid laws:

\[Y_{k,h}=Y_{k+1,0},\qquad B_{k,h}=\sqrt{2h}\,\xi_k,\qquad(Y_{k,0})_\#P=\pi.\]

For \(0\le t\le h\), require

\[Y_{k,t}-Y_{k,0}=-\int_0^t g(Y_{k,s})\,ds+B_{k,t}.\]

Require the score moment bound \(\|g(Y_{k,t})\|_2\le\sqrt{Ld}\) throughout each closed step interval.



Require the increment moment bound \(\|B_{k,t}\|_2\le\sqrt{2d}\sqrt t\) on the same interval.



The segment-joining condition makes the local initial point of segment \(k+1\) equal to the endpoint of segment \(k\). The stationary grid-law condition identifies each comparison marginal with \(\pi\), while the endpoint-noise condition identifies the reference's final increment with the algorithm innovation. The two moment assumptions have different roles: drift controls the deterministic segment increment, and noise controls its square-root growth.

\begin{leanbox}{Lean: \texttt{L2Reference}}
\leansource{\texttt{L2Reference} \quad UnadjustedLangevin/\allowbreak{}HilbertDalalyan.lean, lines 50–65}
\begin{leancode}
\leanline{0}{structure\ L2Reference\ (U\ :\ H\ →\ ℝ)\ (h\ :\ ℝ)\ (L\ :\ ℝ≥0)}
\leanline{4}{(hLip\ :\ LipschitzWith\ L\ (gradient\ U))\ (noise\ :\ ℕ\ →\ Lp\ (H)\ 2\ P)}
\leanline{4}{(π\ :\ Measure\ (H))\ where}
\leanline{2}{segment\ :\ ℕ\ →\ ℝ\ →\ Lp\ (H)\ 2\ P}
\leanline{2}{increment\ :\ ℕ\ →\ ℝ\ →\ Lp\ (H)\ 2\ P}
\leanline{2}{continuous\ :\ ∀\ n,\ Continuous\ (segment\ n)}
\leanline{2}{join\ :\ ∀\ n,\ segment\ n\ h\ =\ segment\ (n\ +\ 1)\ 0}
\leanline{2}{noiseEndpoint\ :\ ∀\ n,\ increment\ n\ h\ =\ Real.sqrt\ (2\ *\ h)\ •\ noise\ n}
\leanline{2}{stationary\ :\ ∀\ n,\ P.map\ (segment\ n\ 0)\ =\ π}
\leanline{2}{integralEquation\ :\ ∀\ n\ t,\ t\ ∈\ Set.Icc\ 0\ h\ →}
\leanline{4}{segment\ n\ t\ -\ segment\ n\ 0\ =}
\leanline{6}{-(∫\ s\ in\ 0..t,\ lpDrift\ P\ (gradient\ U)\ hLip\ (segment\ n\ s))\ +\ increment\ n\ t}
\leanline{2}{scoreBound\ :\ ∀\ n\ t,\ t\ ∈\ Set.Icc\ 0\ h\ →}
\leanline{4}{‖lpDrift\ P\ (gradient\ U)\ hLip\ (segment\ n\ t)‖\ ≤\ Real.sqrt\ ((L\ :\ ℝ)\ *\ d)}
\leanline{2}{noiseBound\ :\ ∀\ n\ t,\ t\ ∈\ Set.Icc\ 0\ h\ →}
\leanline{4}{‖increment\ n\ t‖\ ≤\ Real.sqrt\ (2\ *\ (d\ :\ ℝ))\ *\ Real.sqrt\ t}
\end{leancode}
\leanlabel{How the Lean proof follows the argument}
\texttt{L2\allowbreak{}Reference} collects continuous segments and their additive increments. \texttt{join} identifies adjacent endpoints; \texttt{noise\allowbreak{}Endpoint} matches the increment to the numerical innovation; \texttt{stationary} specifies the law at each segment start. \texttt{integral\allowbreak{}Equation} is the continuous drift equation. The fields \texttt{score\allowbreak{}Bound} and \texttt{noise\allowbreak{}Bound} are bounds in the Hilbert-space norm at exponent two, hence root second-moment bounds. No Gaussian law is asserted by the structure itself.
\end{leanbox}

\section{\texorpdfstring{Convergence from the supplied integral reference}{Convergence from the supplied integral reference}}


Suppose a reference with the preceding properties is supplied, \(h\ge0\), and \(Lh\le2\). Let \(0\le q<1\), and assume the lifted Euler map contracts:

\[\|(X-hg(X))-(Y-hg(Y))\|_2\le q\|X-Y\|_2.\]

For any initial vector \(X_0\in\mathcal H\), the \(L^2\) chain driven by the reference endpoint noise satisfies

\[W_2((X_n)_\#P,\pi)\le q^n\|X_0-Y_{0,0}\|_2+\frac{(33/20)Lh\sqrt{hd}}{1-q}.\]

The statement is conditional on the supplied reference, and the initial term is the cost of the chosen coupling.



Set \(g=\widehat{\nabla U}\), write the numerical iterates as \(X_k\), and define

\[b=\frac{33}{20}Lh\sqrt{hd},\qquad D_k=\int_0^h\bigl(g(Y_{k,t})-g(Y_{k,0})\bigr)\,dt.\]

The bound \(b\) is nonnegative. Apply the local-defect theorem to the gradient map on square-integrable classes and the supplied equation and moment bounds, deriving \(\|D_k\|_2\le b\).



We claim the reference endpoints obey

\[Y_{k+1,0}=Y_{k,0}-hg(Y_{k,0})+\sqrt{2h}\,\xi_k-D_k.\]

The lifted gradient is Lipschitz and the reference segment is continuous, so their composition is continuous and Bochner integrable on \([0,h]\).



Apply the Euler-defect identity at \(t=h\), then substitute the segment join and endpoint noise identities. This proves the claimed reference update. Its defect has a minus sign.



Fix \(k\) and compare the numerical update with the reference endpoint equation, both using the same \(\xi_k\).



Cancel their identical noise terms. Because the reference subtracts \(D_k\),

\[X_{k+1}-Y_{k+1,0}=\bigl(X_k-hg(X_k)\bigr)-\bigl(Y_{k,0}-hg(Y_{k,0})\bigr)+D_k.\]

Apply the triangle inequality, Euler contraction, and the defect bound. Setting \(e_k=\|X_k-Y_{k,0}\|_2\) gives \(e_{k+1}\le qe_k+b\).



The joint law of the pair supplies a finite-cost coupling, so its transport distance is at most \(e_n\). The second marginal is \(\pi\) by the stationary grid-law assumption.



Apply the scalar recurrence bound to the derived inequality and compose it with the transport comparison.



The defect has size \(b=(33/20)Lh\sqrt{hd}\), and this size is uniform in \(k\) because the supplied moment bounds hold on every segment. The recurrence is consequently homogeneous in its constants. Its initial value remains the norm of the supplied pair, so this bound retains the cost of the supplied initial coupling until the subsequent infimum argument.

\begin{leanbox}{Lean: \texttt{l2\_\allowbreak{}reference\_\allowbreak{}bound}}
\leansource{\texttt{l2\_\allowbreak{}reference\_\allowbreak{}bound} \quad UnadjustedLangevin/\allowbreak{}HilbertDalalyan.lean, lines 69–114}
\begin{leancode}
\leanline{0}{theorem\ l2\_reference\_bound}
\leanline{4}{(U\ :\ H\ →\ ℝ)\ (h\ :\ ℝ)\ (L\ :\ ℝ≥0)\ (hLip\ :\ LipschitzWith\ L\ (gradient\ U))}
\leanline{4}{(initial\ :\ Lp\ (H)\ 2\ P)\ (noise\ :\ ℕ\ →\ Lp\ (H)\ 2\ P)\ (π\ :\ Measure\ (H))}
\leanline{4}{(ref\ :\ L2Reference\ P\ d\ U\ h\ L\ hLip\ noise\ π)}
\leanline{4}{(hh\ :\ 0\ ≤\ h)\ (hstep\ :\ (L\ :\ ℝ)\ *\ h\ ≤\ 2)\ (q\ :\ ℝ)\ (hq\ :\ 0\ ≤\ q)\ (hq1\ :\ q\ <\ 1)}
\leanline{4}{(hcontract\ :\ ∀\ X\ Y\ :\ Lp\ (H)\ 2\ P,}
\leanline{6}{‖(X\ -\ h\ •\ lpDrift\ P\ (gradient\ U)\ hLip\ X)\ -}
\leanline{8}{(Y\ -\ h\ •\ lpDrift\ P\ (gradient\ U)\ hLip\ Y)‖\ ≤\ q\ *\ ‖X\ -\ Y‖)\ (n\ :\ ℕ)\ :}
\leanline{4}{wasserstein2Real\ (P.map\ (lpChain\ P\ U\ h\ hLip\ initial\ noise\ n))\ π\ ≤}
\leanline{6}{q\ \textasciicircum{}\ n\ *\ ‖initial\ -\ ref.segment\ 0\ 0‖\ +}
\leanline{8}{((33\ /\ 20\ :\ ℝ)\ *\ (L\ :\ ℝ)\ *\ h\ *\ Real.sqrt\ (h\ *\ d))\ /\ (1\ -\ q)}
\end{leancode}
\leanlabel{How the Lean proof follows the argument}
In \texttt{l2\_\allowbreak{}reference\_\allowbreak{}bound}, \texttt{g} is the lifted drift, \texttt{X} the numerical chain, \texttt{b} the uniform local-error bound and \texttt{D} the segment defect. \texttt{h\allowbreak{}D} applies the integrated local estimate. \texttt{hreference} rewrites the continuous endpoint as an Euler update minus its defect using segment joining and the shared noise endpoint. In \texttt{hrec}, the identity \texttt{hid} cancels the common noise, and the norm triangle inequality gives the scalar recurrence. \texttt{wasserstein2\_\allowbreak{}le\_\allowbreak{}Lp\_\allowbreak{}cost\_\allowbreak{}real} bounds transport by the joint realization, \texttt{ref.\allowbreak{}stationary} identifies the target marginal, and \texttt{affine\_\allowbreak{}recurrence\_\allowbreak{}bound} sums the recurrence.
\end{leanbox}

\section{\texorpdfstring{The small-step bound}{The small-step bound}}


Suppose \(U\) is differentiable and \(m\)-strongly convex on \(H\), \(\nabla U\) is \(L\)-Lipschitz, and \(0<m\le L\). Let \(0<h\le2/(m+L)\), and supply a reference with the preceding properties. For any initial vector \(X_0\in\mathcal H\),

\[W_2((X_n)_\#P,\pi)\le(1-mh)^n\|X_0-Y_{0,0}\|_2+\frac{33}{20}\frac Lm\sqrt{hd}.\]

This uses a chosen initial coupling and a supplied stationary reference. Replacing the initial cost by \(W_2(\mathcal L(X_0),\pi)\) needs the additional uniform family of reference couplings explained in the final section.



Since \(m+L>0\), clear the positive denominator in the step hypothesis to obtain \(h(m+L)\le2\).



Use \(m\le L\) to get \(2mh\le h(m+L)\le2\), hence \(mh\le1\).



The same threshold gives \(Lh\le2\). Thus \(q=1-mh\ge0\); since \(mh>0\), also \(q<1\).



The sharp small-step contraction of the lifted Euler map gives the contraction hypothesis with \(q=1-mh\). Apply the integral-reference theorem.



Because \(m,h>0\), simplify the bias term:

\[\frac{(33/20)Lh\sqrt{hd}}{1-(1-mh)}=\frac{(33/20)Lh\sqrt{hd}}{mh}=\frac{33}{20}\frac Lm\sqrt{hd}.\]

Substitute the bias identity into the integral-reference bound.



The factor \(q=1-mh\) lies in \([0,1)\) under the positive-curvature small-step assumptions. Its gap is exactly \(mh\), so dividing the local \(h^{3/2}\) error by that gap produces the global \(\sqrt h\) bias. The remaining factor \(L/m\) is the condition-number dependence of this estimate.

\begin{leanbox}{Lean: \texttt{dalalyan\_\allowbreak{}small\_\allowbreak{}step\_\allowbreak{}of\_\allowbreak{}integral\_\allowbreak{}reference}}
\leansource{\texttt{dalalyan\_\allowbreak{}small\_\allowbreak{}step\_\allowbreak{}of\_\allowbreak{}integral\_\allowbreak{}reference} \quad UnadjustedLangevin/\allowbreak{}HilbertDalalyan.lean, lines 118–143}
\begin{leancode}
\leanline{0}{theorem\ dalalyan\_small\_step\_of\_integral\_reference}
\leanline{4}{(U\ :\ H\ →\ ℝ)\ (m\ h\ :\ ℝ)\ (L\ :\ ℝ≥0)}
\leanline{4}{(hdiff\ :\ Differentiable\ ℝ\ U)\ (hsc\ :\ StrongConvexOn\ Set.univ\ m\ U)}
\leanline{4}{(hLip\ :\ LipschitzWith\ L\ (gradient\ U))\ (hm\ :\ 0\ <\ m)\ (hmL\ :\ m\ ≤\ (L\ :\ ℝ))}
\leanline{4}{(hh\ :\ 0\ <\ h)\ (hsmall\ :\ h\ ≤\ 2\ /\ (m\ +\ (L\ :\ ℝ)))}
\leanline{4}{(initial\ :\ Lp\ (H)\ 2\ P)\ (noise\ :\ ℕ\ →\ Lp\ (H)\ 2\ P)\ (π\ :\ Measure\ (H))}
\leanline{4}{(ref\ :\ L2Reference\ P\ d\ U\ h\ L\ hLip\ noise\ π)\ (n\ :\ ℕ)\ :}
\leanline{4}{wasserstein2Real\ (P.map\ (lpChain\ P\ U\ h\ hLip\ initial\ noise\ n))\ π\ ≤}
\leanline{6}{(1\ -\ m\ *\ h)\ \textasciicircum{}\ n\ *\ ‖initial\ -\ ref.segment\ 0\ 0‖\ +}
\leanline{8}{(33\ /\ 20\ :\ ℝ)\ *\ ((L\ :\ ℝ)\ /\ m)\ *\ Real.sqrt\ (h\ *\ d)}
\end{leancode}
\leanlabel{How the Lean proof follows the argument}
\texttt{dalalyan\_\allowbreak{}small\_\allowbreak{}step\_\allowbreak{}of\_\allowbreak{}integral\_\allowbreak{}reference} receives \texttt{ref}, an \texttt{L2\allowbreak{}Reference} whose fields contain the segment equations, target laws and moment bounds. \texttt{hproduct}, \texttt{hmh} and \texttt{hstep} derive the scalar step consequences. These give the nonnegative contraction factor \texttt{hq} and its strict upper bound \texttt{hq1}. The invocation of \texttt{l2\_\allowbreak{}reference\_\allowbreak{}bound} combines the integral local defect with \texttt{lp\allowbreak{}Drift\_\allowbreak{}gradient\_\allowbreak{}contract}. \texttt{hid} cancels the positive step from the forcing divided by the contraction gap, producing the displayed constant. This theorem is conditional on its reference; the preceding stationarity and product-construction results supply that reference for the final Gaussian algorithm.
\end{leanbox}

\section{\texorpdfstring{The large-step bound}{The large-step bound}}


Suppose \(U\) is differentiable and \(m\)-strongly convex on \(H\), \(\nabla U\) is \(L\)-Lipschitz, and \(0<m\le L\). Let \(h>0\), \(h\ge2/(m+L)\), and \(Lh<2\), and supply a reference with the stated properties. For any initial vector \(X_0\in\mathcal H\),

\[W_2((X_n)_\#P,\pi)\le(Lh-1)^n\|X_0-Y_{0,0}\|_2+\frac{33}{20}\frac{Lh}{2-Lh}\sqrt{hd}.\]

The denominator is positive by strict stability. The reference remains an input, and the initial term remains the chosen coupling cost.



Since \(m+L>0\), clear the denominator in the lower step threshold to obtain \(2\le h(m+L)\).



Use \(m\le L\) and \(h\ge0\): \(2\le h(m+L)\le2Lh\). Thus \(Lh\ge1\) and \(q=Lh-1\ge0\).



Strict stability \(Lh<2\) gives \(q=Lh-1<1\).



Apply the sharp large-step Euler contraction with \(q=Lh-1\), then invoke the integral-reference theorem.



Use \(1-(Lh-1)=2-Lh\) and factor out \(\sqrt{hd}\):

\[\frac{(33/20)Lh\sqrt{hd}}{1-(Lh-1)}=\left(\frac{(33/20)Lh}{2-Lh}\right)\sqrt{hd}.\]

Substitute this identity into the integral-reference bound to obtain the large-step formula.



Here \(q=Lh-1\) and the gap is \(1-q=2-Lh>0\). Dividing the same local-defect estimate by that gap gives the displayed additive large-step bound. As the denominator decreases to zero, this geometric-sum upper bound increases. The strict condition on the step keeps the denominator positive.

\begin{leanbox}{Lean: \texttt{dalalyan\_\allowbreak{}large\_\allowbreak{}step\_\allowbreak{}of\_\allowbreak{}integral\_\allowbreak{}reference}}
\leansource{\texttt{dalalyan\_\allowbreak{}large\_\allowbreak{}step\_\allowbreak{}of\_\allowbreak{}integral\_\allowbreak{}reference} \quad UnadjustedLangevin/\allowbreak{}HilbertDalalyan.lean, lines 146–168}
\begin{leancode}
\leanline{0}{theorem\ dalalyan\_large\_step\_of\_integral\_reference}
\leanline{4}{(U\ :\ H\ →\ ℝ)\ (m\ h\ :\ ℝ)\ (L\ :\ ℝ≥0)}
\leanline{4}{(hdiff\ :\ Differentiable\ ℝ\ U)\ (hsc\ :\ StrongConvexOn\ Set.univ\ m\ U)}
\leanline{4}{(hLip\ :\ LipschitzWith\ L\ (gradient\ U))\ (hm\ :\ 0\ <\ m)\ (hmL\ :\ m\ ≤\ (L\ :\ ℝ))}
\leanline{4}{(hh\ :\ 0\ <\ h)\ (hlarge\ :\ 2\ /\ (m\ +\ (L\ :\ ℝ))\ ≤\ h)\ (hstep\ :\ (L\ :\ ℝ)\ *\ h\ <\ 2)}
\leanline{4}{(initial\ :\ Lp\ (H)\ 2\ P)\ (noise\ :\ ℕ\ →\ Lp\ (H)\ 2\ P)\ (π\ :\ Measure\ (H))}
\leanline{4}{(ref\ :\ L2Reference\ P\ d\ U\ h\ L\ hLip\ noise\ π)\ (n\ :\ ℕ)\ :}
\leanline{4}{wasserstein2Real\ (P.map\ (lpChain\ P\ U\ h\ hLip\ initial\ noise\ n))\ π\ ≤}
\leanline{6}{((L\ :\ ℝ)\ *\ h\ -\ 1)\ \textasciicircum{}\ n\ *\ ‖initial\ -\ ref.segment\ 0\ 0‖\ +}
\leanline{8}{((33\ /\ 20\ :\ ℝ)\ *\ (L\ :\ ℝ)\ *\ h\ /\ (2\ -\ (L\ :\ ℝ)\ *\ h))\ *\ Real.sqrt\ (h\ *\ d)}
\end{leancode}
\leanlabel{How the Lean proof follows the argument}
\texttt{dalalyan\_\allowbreak{}large\_\allowbreak{}step\_\allowbreak{}of\_\allowbreak{}integral\_\allowbreak{}reference} derives \texttt{hproduct} from the lower step threshold. Together with the order of curvature and smoothness, this proves \texttt{hq}, nonnegativity of the large-step factor; \texttt{hstep} gives \texttt{hq1}, its strict upper bound. The proof invokes \texttt{l2\_\allowbreak{}reference\_\allowbreak{}bound} with \texttt{lp\allowbreak{}Drift\_\allowbreak{}gradient\_\allowbreak{}contract\_\allowbreak{}large}. The identity \texttt{hid} rewrites the contraction gap as the positive large-step denominator and moves the square-root factor outside the quotient.
\end{leanbox}

\end{document}
